\documentclass[11pt]{article}
\usepackage[a4paper,margin=22mm]{geometry}
\usepackage{fontspec}
\setmainfont{DejaVu Serif}
\setsansfont{DejaVu Sans}
\usepackage{microtype}
\usepackage{xcolor}
\usepackage{hyperref}
\usepackage{enumitem}
\usepackage{parskip}
\definecolor{Accent}{HTML}{971D32}
\definecolor{Muted}{HTML}{666666}
\hypersetup{colorlinks=true,urlcolor=Accent,linkcolor=Accent}
\setlist{leftmargin=1.6em,itemsep=0.3em,topsep=0.35em}
\setcounter{tocdepth}{2}
\renewcommand{\contentsname}{Contents}
\newcommand{\sectionrule}{\par\vspace{-0.5em}\noindent\textcolor{Accent}{\rule{\linewidth}{0.6pt}}\par}
\begin{document}

\begin{center}
{\sffamily\bfseries\color{Accent} TOEFL WORD OF THE DAY\par}
\vspace{8mm}
{\fontsize{42}{48}\selectfont\bfseries democratize\par}
\vspace{2mm}
{\Large\sffamily\color{Accent} /dɪˈmɑːkrəˌtaɪz/\par}
\vspace{2mm}
{\small\sffamily Local audio phrase: democratize\par}
{\small\sffamily Audio file: audios/democratize.mp3\par}
\vspace{2mm}
{\large\sffamily\color{Muted} transitive verb\par}
\vspace{5mm}
{\sffamily make a system more democratic \textbullet\ make something accessible to more people\par}
\vspace{6mm}
{\large To democratize something is to spread power, participation, or access beyond a small group.\par}
\vspace{6mm}
\begin{minipage}{0.84\textwidth}
\itshape “Online libraries can democratize access to academic research.”
\end{minipage}\par
\vspace{10mm}
\textbf{Date:} July 27th 2026\quad\textbullet\quad \textbf{Level:} B1--mid-B2 TOEFL
\vfill
Created and compiled by \textbf{Amir Kooshky}\par
\vspace{2mm}
\href{https://t.me/KooshkyTOEFL}{Telegram Channel}\quad\textbullet\quad
\href{https://www.instagram.com/kooshkytoefl}{Instagram}\quad\textbullet\quad
\href{https://t.me/KooshkyTOEFL\_pv}{Personal Telegram}
\end{center}
\newpage
\tableofcontents
\newpage

\section{Meanings and usage}
\sectionrule
\subsection{Meaning 1: Make more democratic}
\textbf{Part of speech:} transitive verb

\textbf{IPA:} /dɪˈmɑːkrəˌtaɪz/

\textbf{Pronunciation phrase:} democratize

\textbf{Local audio file:} audios/democratize.mp3

\textbf{Register:} formal; common in politics, history, sociology, and discussions of organizations

\textbf{Definition:} to make a country, organization, or system more democratic by giving people more power or a greater role in decisions

\textbf{In simpler words:} give more people a real voice in decisions

\textbf{Typical pattern:} democratize + country, institution, system, or decision-making

\subsubsection{Natural examples}
\begin{enumerate}
  \item The reforms were intended to democratize the country's political system.
  \item Student representatives argued that the university should democratize its decision-making process.
  \item The new constitution helped democratize several public institutions.
  \item Local councils were created to democratize the management of community resources.
  \item The movement sought to democratize the workplace by giving employees a stronger voice.
  \item Critics questioned whether the changes would truly democratize the organization.
\end{enumerate}

\subsubsection{Nuance and implication}
\textbf{Central nuance:} The word suggests moving power or decision-making away from a small group and toward wider participation.

\textbf{Usual implication:} It often describes a process rather than a single immediate change.

\textbf{Compared with liberalize:} To liberalize is to reduce restrictions or government control. To democratize is specifically to increase democratic participation, representation, or accountability.

\subsubsection{Grammar notes}
\textbf{How to use it:} This verb normally needs the country, organization, system, or process that is being changed after it.

\textbf{What can follow it:} It is often followed by words such as government, society, institution, system, workplace, decision-making, or governance.

\textbf{Position in a sentence:} It is also common in the passive form, as in The institution was gradually democratized.

\subsubsection{Useful patterns}
\textbf{Pattern:} democratize + institution or system

\textbf{Use:} Use this when reforms give more people influence over how an organization or system operates.

\textbf{Example:} The proposal aims to democratize the national education system.

\textbf{Pattern:} democratize decision-making

\textbf{Use:} Use this when a wider group is given a meaningful role in making choices.

\textbf{Example:} The committee introduced open meetings to democratize decision-making.

\textbf{Pattern:} be democratized

\textbf{Use:} Use the passive form when the change itself is more important than the person or group causing it.

\textbf{Example:} The organization was gradually democratized through a series of internal reforms.

\subsubsection{Common wrong usage}
\paragraph{Correction 1}
\textbf{Wrong:} The reform democratized people to make decisions.

\textbf{Correct:} The reform democratized the decision-making process.

\textbf{Why:} Democratize is normally followed by the system or process that becomes more democratic, not by people plus to.

\paragraph{Correction 2}
\textbf{Wrong:} The government wants to democratize in the public institutions.

\textbf{Correct:} The government wants to democratize the public institutions.

\textbf{Why:} In this meaning, the verb normally takes its object directly, without in.

\subsection{Meaning 2: Widen access}
\textbf{Part of speech:} transitive verb

\textbf{IPA:} /dɪˈmɑːkrəˌtaɪz/

\textbf{Pronunciation phrase:} democratize

\textbf{Local audio file:} audios/democratize.mp3

\textbf{Register:} formal and increasingly common in education, technology, business, and public policy

\textbf{Definition:} to make information, education, technology, or another useful resource available to many more people instead of only a small or powerful group

\textbf{In simpler words:} make something available to far more people

\textbf{Typical pattern:} democratize access to + resource or democratize + knowledge, education, technology, or opportunity

\subsubsection{Natural examples}
\begin{enumerate}
  \item Online libraries can democratize access to academic research.
  \item Low-cost smartphones have helped democratize access to digital services.
  \item The museum hopes to democratize art by offering free virtual exhibitions.
  \item Open-source software can democratize advanced technology for small organizations.
  \item Public data projects aim to democratize information that was once difficult to obtain.
  \item Some educators believe artificial intelligence could democratize personalized learning.
\end{enumerate}

\subsubsection{Nuance and implication}
\textbf{Central nuance:} This use is not necessarily about elections or government. It describes removing barriers that kept a valuable resource within a limited group.

\textbf{Usual implication:} The barriers may involve cost, location, technical knowledge, social position, or institutional control.

\textbf{Compared with popularize:} To popularize something is to make it widely known or liked. To democratize it is to make access or participation more widely available.

\subsubsection{Grammar notes}
\textbf{How to use it:} This verb normally needs the resource, opportunity, or type of access that is being made more widely available after it.

\textbf{What can follow it:} Common objects include access, education, knowledge, information, technology, finance, culture, and opportunity.

\textbf{Common preposition:} A very common pattern is democratize access to something, as in democratize access to higher education.

\textbf{Position in a sentence:} The phrase help democratize is common when one development contributes to wider access but is not the only cause.

\subsubsection{Useful patterns}
\textbf{Pattern:} democratize access to + noun

\textbf{Use:} Use this when more people are given a realistic chance to obtain or use something.

\textbf{Example:} Digital courses may democratize access to professional training.

\textbf{Pattern:} democratize + knowledge, technology, or education

\textbf{Use:} Use this when the resource itself becomes less limited to experts, wealthy users, or major institutions.

\textbf{Example:} The platform was designed to democratize financial knowledge.

\textbf{Pattern:} help democratize + noun

\textbf{Use:} Use this when something contributes to broader access.

\textbf{Example:} Affordable equipment could help democratize scientific research.

\subsubsection{Common wrong usage}
\paragraph{Correction 1}
\textbf{Wrong:} The website democratizes the access to legal information.

\textbf{Correct:} The website democratizes access to legal information.

\textbf{Why:} Access is normally used without the in this pattern.

\paragraph{Correction 2}
\textbf{Wrong:} The program democratized education to rural students.

\textbf{Correct:} The program democratized access to education for rural students.

\textbf{Why:} Use access to for the resource and for to identify the people who benefit.

\section{Word family}
\sectionrule
\subsection{democratization}
\textbf{Part of speech:} uncountable noun

\textbf{IPA:} /dɪˌmɑːkrətəˈzeɪʃən/

\textbf{Definition:} the process of making a country, organization, system, or resource more democratic or more widely accessible

\textbf{Relationship to the headword:} This noun names the process or result of democratizing something.

\textbf{Grammar pattern:} the democratization of + country, institution, system, knowledge, or access

\textbf{Usage note:} It is usually used without a or an.

\subsubsection{Examples}
\begin{enumerate}
  \item The report examines the democratization of local government.
  \item Digital tools have contributed to the democratization of knowledge.
\end{enumerate}

\subsection{democratic}
\textbf{Part of speech:} adjective

\textbf{IPA:} /ˌdɛməˈkrætɪk/

\textbf{Definition:} based on democracy or giving people a fair and meaningful role in decisions

\textbf{Relationship to the headword:} This adjective describes a system, process, or organization that has democratic qualities.

\textbf{Grammar pattern:} democratic + government, election, institution, process, or society

\subsubsection{Examples}
\begin{enumerate}
  \item The country held its first democratic election in decades.
  \item Employees called for a more democratic workplace.
\end{enumerate}

\subsection{democracy}
\textbf{Part of speech:} countable or uncountable noun

\textbf{IPA:} /dɪˈmɑːkrəsi/

\textbf{Definition:} a system of government in which people choose their leaders and can take part in political decisions

\textbf{Relationship to the headword:} Democratize originally refers to making something more like a democracy.

\textbf{Grammar pattern:} democracy or a democracy

\textbf{Usage note:} Use democracy without a number for the general idea and a democracy for a country or political system.

\subsubsection{Examples}
\begin{enumerate}
  \item Free elections are an important part of democracy.
  \item The transition to a stable democracy took many years.
\end{enumerate}

\section{Common collocations}
\sectionrule
\subsection{democratize access to education}
\textbf{Applies to:} Widen access

\textbf{Meaning:} make education realistically available to a wider group

\textbf{Pattern:} democratize access to + education or training

\textbf{Example:} Scholarship programs can help democratize access to education.

\subsection{democratize access to information}
\textbf{Applies to:} Widen access

\textbf{Meaning:} allow far more people to obtain information

\textbf{Pattern:} democratize access to + information or data

\textbf{Example:} Public databases have democratized access to scientific information.

\subsection{democratize technology}
\textbf{Applies to:} Widen access

\textbf{Meaning:} make technology usable or affordable for a wider population

\textbf{Example:} Affordable tools can democratize technology that was once limited to large companies.

\subsection{democratize knowledge}
\textbf{Applies to:} Widen access

\textbf{Meaning:} make knowledge available beyond a small expert or privileged group

\textbf{Example:} The project aims to democratize knowledge by publishing its courses for free.

\subsection{democratize decision-making}
\textbf{Applies to:} Make more democratic

\textbf{Meaning:} give a wider group a meaningful role in making decisions

\textbf{Example:} The reform was designed to democratize decision-making within the organization.

\subsection{democratize an institution}
\textbf{Applies to:} Make more democratic

\textbf{Meaning:} change an institution so that power and participation are more widely shared

\textbf{Example:} Members voted for reforms intended to democratize the institution.

\subsection{democratize a political system}
\textbf{Applies to:} Make more democratic

\textbf{Meaning:} make a political system more representative and accountable

\textbf{Example:} The constitution introduced measures to democratize the political system.

\subsection{the democratization of knowledge}
\textbf{Applies to:} Widen access

\textbf{Meaning:} the process of making knowledge available to more people

\textbf{Example:} The internet has contributed to the democratization of knowledge.

\textbf{Usage note:} Democratization is the noun form.

\subsection{the democratization of society}
\textbf{Applies to:} Make more democratic

\textbf{Meaning:} the process through which social and political power becomes more widely shared

\textbf{Example:} The study traces the gradual democratization of society after the conflict.

\textbf{Usage note:} Democratization is the noun form.

\section{Synonyms and antonyms}
\sectionrule
\subsection{Close synonyms}
\subsubsection{liberalize}
\textbf{Part of speech:} transitive verb

\textbf{Applies to:} Make more democratic

\textbf{Shared meaning:} Both words can describe reforms that make a system more open.

\textbf{Difference:} Liberalize means to reduce restrictions or controls. Such reforms may support democratization, but they do not necessarily give people more political power.

\textbf{Example:} The government liberalized restrictions on the press.

\subsubsection{broaden access to}
\textbf{Part of speech:} verb phrase

\textbf{Applies to:} Widen access

\textbf{Shared meaning:} Both expressions mean making a resource available to more people.

\textbf{Difference:} Broaden access to is a more direct and neutral phrase. Democratize often adds the idea that access was previously controlled by a privileged or powerful group.

\textbf{Example:} The program broadened access to university courses.

\subsubsection{popularize}
\textbf{Part of speech:} transitive verb

\textbf{Applies to:} Widen access

\textbf{Shared meaning:} Both can describe something reaching a much wider audience.

\textbf{Difference:} Popularize means to make something widely known, understood, or liked. Democratize focuses on removing barriers to access or participation.

\textbf{Example:} The television series popularized astronomy among young viewers.

\section{Etymology}
\sectionrule
\textbf{Origin:} Formed from democrat and the verb-forming suffix -ize, with the basic meaning of making something democratic.

\section{Practice exercises}
\sectionrule
\subsection{Word family choice}
Choose the best member of the word family for each blank.

\begin{enumerate}
  \item The report describes the gradual \_\_\_\_\_ of the country's political institutions.
  \item The reform created a more \_\_\_\_\_ system of local government.
  \item Free elections alone do not always create a stable \_\_\_\_\_.
\end{enumerate}

\subsection{Correct or incorrect?}
Decide whether each sentence is correct. Correct the inaccurate ones.

\begin{enumerate}
  \item The platform aims to democratize access to legal advice.
  \item The reforms democratized in the national parliament.
  \item The website democratized the access to historical documents.
  \item The new voting rules were introduced to democratize decision-making.
\end{enumerate}

\subsection{Collocation practice}
Complete each sentence with a natural collocation.

\begin{enumerate}
  \item Free online courses may help democratize \_\_\_\_\_ to higher education.
  \item The council introduced public meetings to democratize \_\_\_\_\_.
  \item The project seeks to democratize \_\_\_\_\_ by making research articles freely available.
\end{enumerate}

\subsection{Complete answer key}
\subsubsection{Word family choice}
\begin{enumerate}
  \item \textbf{democratization.} The report describes the gradual democratization of the country's political institutions. \emph{Use the noun democratization after the.}
  \item \textbf{democratic.} The reform created a more democratic system of local government. \emph{Use the adjective democratic before the noun system.}
  \item \textbf{democracy.} Free elections alone do not always create a stable democracy. \emph{Use the countable noun democracy for a particular political system or country.}
\end{enumerate}

\subsubsection{Correct or incorrect?}
\begin{enumerate}
  \item \textbf{Correct} \emph{Democratize access to is a natural pattern for making a resource available to more people.}
  \item \textbf{Incorrect}. The reforms democratized the national parliament. \emph{The verb normally takes the institution or system directly as its object.}
  \item \textbf{Incorrect}. The website democratized access to historical documents. \emph{Access is normally used without the in this expression.}
  \item \textbf{Correct} \emph{Democratize decision-making correctly means giving more people a real role in decisions.}
\end{enumerate}

\subsubsection{Collocation practice}
\begin{enumerate}
  \item \textbf{access.} Free online courses may help democratize access to higher education. \emph{Democratize access to is an established pattern.}
  \item \textbf{decision-making.} The council introduced public meetings to democratize decision-making. \emph{Democratize decision-making means giving a wider group influence over choices.}
  \item \textbf{knowledge.} The project seeks to democratize knowledge by making research articles freely available. \emph{Democratize knowledge means making knowledge available beyond a limited group.}
\end{enumerate}

\vfill
\begin{center}
\small\color{Muted} Created and compiled by \textbf{Amir Kooshky}\\
\href{https://t.me/KooshkyTOEFL}{Telegram Channel}\quad\textbullet\quad
\href{https://www.instagram.com/kooshkytoefl}{Instagram}\quad\textbullet\quad
\href{https://t.me/KooshkyTOEFL\_pv}{Personal Telegram}
\end{center}
\end{document}
